% Luc Destombe --- licence CC BY-NC-SA 4.0
% https://creativecommons.org/licenses/by-nc-sa/4.0/deed.fr
% Fichier autonome : compiler avec pdflatex, deux fois (signets, nombre de pages).
\documentclass[11pt,a4paper]{article}
\makeatletter
\newif\iflycee@unecolonne
\lycee@unecolonnetrue
\newif\iflycee@palatino
\lycee@palatinotrue

\RequirePackage[utf8]{inputenc}
\RequirePackage[T1]{fontenc}
\RequirePackage[french]{babel}
\RequirePackage{etoolbox}

\newcommand{\ShorthandsOff}{\@ifundefined{shorthandoff}{}{\shorthandoff{;:!?}}}
\newcommand{\ShorthandsOn}{\@ifundefined{shorthandon}{}{\shorthandon{;:!?}}}
\ShorthandsOff
\AtBeginDocument{\ShorthandsOn}
\AtBeginEnvironment{tikzpicture}{\ShorthandsOff}

\RequirePackage{geometry}
\RequirePackage{fancyhdr}
\RequirePackage{lastpage}
\RequirePackage{multicol}
\RequirePackage{array}
\RequirePackage{enumitem}
\RequirePackage{xcolor}
\RequirePackage{needspace}

\RequirePackage{amsmath,amssymb}
\RequirePackage{mathpazo}
\DeclareMathSizes{10.95}{10.95}{8}{5.5}
\begingroup\catcode`\_=\active \catcode`\^=\active
\gdef_#1{\sb{\scriptscriptstyle #1}}
\gdef^#1{\sp{\scriptscriptstyle\lycee@moinsepais #1}}
\endgroup
\newcommand\lycee@moins{\mathbin{%
  \setbox\z@\hbox{$\m@th\scriptscriptstyle\lycee@moinsorig$}%
  \dimen@=.55\wd\z@ \advance\dimen@-.5pt
  \hbox to.75\wd\z@{\hss$\m@th\scriptscriptstyle\vcenter{\hbox{%
    \tikz\draw[line width=.5pt,line cap=round](0,0)--(\the\dimen@,0);}}$\hss}}}
\begingroup\lccode`\~=`\- \lowercase{\endgroup\let~}\lycee@moins
\newcommand\lycee@moinsepais{\mathcode`\-="8000 }
\AtBeginDocument{\mathchardef\lycee@moinsorig=\mathcode`\-\relax}
\AtBeginDocument{\catcode`\_=12 \mathcode`\_="8000
  \catcode`\^=12 \mathcode`\^="8000 }
\newcommand\lycee@exposants{%
  \fontdimen13\textfont2=.48\fontdimen6\textfont2
  \fontdimen14\textfont2=.48\fontdimen6\textfont2
  \fontdimen15\textfont2=.48\fontdimen6\textfont2
  \fontdimen13\scriptfont2=.48\fontdimen6\scriptfont2
  \fontdimen14\scriptfont2=.48\fontdimen6\scriptfont2
  \fontdimen15\scriptfont2=.48\fontdimen6\scriptfont2}
\every@math@size\expandafter{\the\every@math@size\lycee@exposants}
\newlength{\retraitcalculs}
\setlength{\retraitcalculs}{2em}

\RequirePackage{tikz}
\usetikzlibrary{arrows.meta,calc,babel,shapes.misc}
\RequirePackage[most]{tcolorbox}
\tcbuselibrary{skins,breakable}

\definecolor{coulDef}{HTML}{1F4E79}
\definecolor{coulProp}{HTML}{7030A0}
\definecolor{coulRem}{HTML}{C55A11}
\definecolor{coulEx}{HTML}{2E7D32}
\definecolor{coulExo}{HTML}{B03A2E}
\definecolor{coulPart}{HTML}{1F4E79}
\definecolor{coulMeth}{HTML}{1B6E4F}
\definecolor{coulCourbe}{HTML}{0B5ED7}
\definecolor{coulCourbeB}{HTML}{C00000}
\definecolor{coulCourbeC}{HTML}{2E7D32}
\definecolor{coulGrille}{HTML}{8C8C8C}
\definecolor{coulRep}{HTML}{1B6E4F}
\definecolor{coulBar}{HTML}{2E7D32}

\newcommand{\R}{\mathbb{R}}
\newcommand{\Rs}{\mathbb{R}^{*}}
\renewcommand{\le}{\leqslant}
\renewcommand{\ge}{\geqslant}
\renewcommand{\approx}{\simeq}
\newcommand{\vect}[1]{\overrightarrow{#1}}
\newcommand{\norme}[1]{\left\lVert #1 \right\rVert}
\newcommand{\intervalle}[2]{\left[#1\,;#2\right]}
\RequirePackage{cancel}
\renewcommand{\CancelColor}{\color{coulExo}}
\newcommand{\fois}[1]{\times{\color{coulExo}#1}}
\newcommand{\barre}[1]{\cancel{#1}}

\tikzset{grille/.style={coulGrille,line width=0.4pt}}
\newcommand{\quadrillage}[4]{\draw[grille,step=1] (#1,#2) grid (#3,#4);}
\tikzset{
  croix/.style={cross out,draw,line width=0.6pt,inner sep=0pt,minimum size=4pt},
  croixdroite/.style={croix,rotate=45,line width=0.8pt},
}
\newcommand{\pt}[3]{\path (#1) node[croix]{} node[#2,font=\small,inner sep=2.5pt]{$#3$};}

\newcommand{\lycee@repere}[4]{%
  \draw[grille,step=1] (#1,#2) grid (#3,#4);
  \draw[-{Stealth[length=2mm]},thick] (#1,0)--({#3+0.4},0) node[below left=-1pt]{$x$};
  \draw[-{Stealth[length=2mm]},thick] (0,#2)--(0,{#4+0.4}) node[below left=-1pt]{$y$};
  \node[below left,font=\scriptsize,inner sep=1.5pt] at (0,0) {$0$};}
\newcommand{\lycee@gradx}[1]{\foreach \k/\l in {#1}{%
  \draw (\k,0.07)--(\k,-0.07) node[below,font=\scriptsize,inner sep=1.5pt]{$\l$};}}
\newcommand{\lycee@grady}[1]{\foreach \k/\l in {#1}{%
  \draw (0.07,\k)--(-0.07,\k) node[left,font=\scriptsize,inner sep=1.5pt]{$\l$};}}
\newcommand{\lycee@intff}[2]{\left[#1\,;#2\right]}
\newcommand{\lycee@intfo}[2]{\left[#1\,;#2\right[}
\newcommand{\lycee@intof}[2]{\left]#1\,;#2\right]}
\newcommand{\lycee@intoo}[2]{\left]#1\,;#2\right[}
\newcommand{\lycee@ens}[1]{\left\{#1\right\}}
\AtBeginDocument{%
  \providecommand{\repere}{\lycee@repere}%
  \providecommand{\gradx}{\lycee@gradx}%
  \providecommand{\grady}{\lycee@grady}%
  \providecommand{\Cf}{\mathcal{C}\sb{\scriptscriptstyle f}}%
  \providecommand{\Cg}{\mathcal{C}\sb{\scriptscriptstyle g}}%
  \providecommand{\intff}{\lycee@intff}%
  \providecommand{\intfo}{\lycee@intfo}%
  \providecommand{\intof}{\lycee@intof}%
  \providecommand{\intoo}{\lycee@intoo}%
  \providecommand{\ens}{\lycee@ens}}

\newlist{questions}{enumerate}{3}
\setlist[questions,1]{label=\arabic*\textdegree),leftmargin=*,itemsep=2pt,topsep=3pt}
\setlist[questions,2]{label=\alph*),leftmargin=*,itemsep=1pt,topsep=2pt}
\setlist[questions,3]{label=\roman*),leftmargin=*,itemsep=1pt,topsep=1pt}
\setlist[itemize]{label=\textbullet,leftmargin=*,itemsep=2pt,topsep=3pt}

\newcommand{\besoinplace}[1]{\par\penalty\@M\begingroup
  \dimen@=#1\relax\dimen@ii=\pagegoal\advance\dimen@ii-\pagetotal
  \ifdim\dimen@>\dimen@ii\ifdim\dimen@ii>\z@\vfil\fi\break\fi\endgroup}

\newcommand{\lycee@pied}{\thepage/\pageref{LastPage}}
\newcommand{\entete}[2]{%
  \pagestyle{fancy}\fancyhf{}%
  \fancyhead[L]{\footnotesize\color{coulPart}\textbf{#1}}%
  \fancyhead[R]{\footnotesize\color{coulPart}\textbf{#2}}%
  \fancyfoot[C]{\footnotesize\lycee@pied}%
  \renewcommand{\headrulewidth}{0.4pt}%
  \renewcommand{\headrule}{\hbox to\headwidth{%
    \color{coulPart!40}\leaders\hrule height \headrulewidth\hfill}}}

\RequirePackage{ccicons}
\newcommand{\lycee@licence}{{\scriptsize\color{gray}%
  \copyright\ Luc Destombe --- \ccbyncsa\ CC BY-NC-SA 4.0}}
\newif\iflycee@licence \lycee@licencetrue
\newcommand{\sanslicence}{\lycee@licencefalse}
\AtBeginDocument{\iflycee@licence\AddToHookNext{shipout/foreground}{%
  \put(0,0){\rlap{\kern\dimexpr1in+\hoffset+\oddsidemargin+\textwidth\relax
    \llap{\raisebox{-\dimexpr1in+\voffset+\topmargin+\headheight+\headsep
      +\textheight+\footskip\relax}{\lycee@licence}}}}}\fi}

\newcommand{\formulesserrees}{%
  \setlength{\abovedisplayskip}{3pt}\setlength{\belowdisplayskip}{3pt}%
  \setlength{\abovedisplayshortskip}{2pt}\setlength{\belowdisplayshortskip}{3pt}}

\setlength{\parindent}{0pt}

\RequirePackage[implicit=false,hidelinks,pdfencoding=auto,psdextra,bookmarksopen,
  bookmarksopenlevel=1,pdfcreator={},pdfproducer={}]{hyperref}
\RequirePackage{bookmark}
\hypersetup{pdfauthor={Luc Destombe},pdfkeywords={CC BY-NC-SA 4.0}}
\newcounter{lycee@signet}
\newcommand{\signet}[3][]{\stepcounter{lycee@signet}%
  \edef\lycee@dest{\if\relax\detokenize{#1}\relax
    signet.\arabic{lycee@signet}\else#1\fi}%
  \hypertarget{\lycee@dest}{}\bookmark[level=#2,dest=\lycee@dest]{#3}}
\pdfstringdefDisableCommands{%
  \def\R{R}\def\vect#1{#1}\def\Cf{Cf}\def\Cg{Cg}\def\fois#1{×#1}%
  \def\barre#1{#1}\def\intervalle#1#2{[#1;#2]}\def\intff#1#2{[#1;#2]}%
  \def\textsuperscript#1{#1}\def\dfrac#1#2{#1/#2}\def\frac#1#2{#1/#2}%
  \def\vec#1{#1⃗}}%
\RequirePackage[official]{eurosym}

\tcbset{exercice corrige/.style={
  enhanced,breakable,before upper=\raggedright,
  colback=white,colframe=coulExo,
  boxrule=0.8pt,arc=2pt,
  left=6pt,right=6pt,top=8pt,bottom=5pt,
  before skip=9pt,after skip=4pt,
  coltitle=white,fonttitle=\bfseries\small,
  attach boxed title to top left={xshift=8pt,yshift=-\tcboxedtitleheight/2},
  boxed title style={colback=coulExo,arc=2pt,boxrule=0pt}}}
\newcounter{exo}
\newtcolorbox{exercice}[1][]{exercice corrige,
  phantom={\signet[exercice-\arabic{exo}]{2}{Exercice \arabic{exo}}},
  title={Exercice \theexo\ifx\relax#1\relax\else{} --- #1\fi}}
\newenvironment{exo}[1][\relax]{\refstepcounter{exo}\begin{exercice}[#1]}{\end{exercice}}

\geometry{top=1.8cm,bottom=1cm,left=1cm,right=1cm,headheight=14pt,footskip=0.6cm}
\renewcommand{\lycee@pied}{\thepage}

\renewcommand{\theexo}{\ifnum\value{exo}<10 0\fi\arabic{exo}}
\tcbset{exercice corrige/.style={
  enhanced,breakable,before upper=\raggedright,
  colback=white,colframe=coulExo,
  boxrule=0.9pt,arc=2pt,
  left=8pt,right=8pt,top=8pt,bottom=6pt,
  before skip=8pt,after skip=6pt,
  coltitle=white,fonttitle=\bfseries,
  attach boxed title to top left={xshift=8pt,yshift=-\tcboxedtitleheight/2},
  boxed title style={colback=coulExo,arc=2pt,boxrule=0pt}}}

\newcommand{\entetefiche}[2]{%
  \begin{center}
    \begin{tcolorbox}[enhanced,width=0.92\linewidth,colback=white,colframe=coulPart,
        boxrule=1.4pt,arc=3pt,halign=center,top=6pt,bottom=6pt]
      {\large\bfseries\color{coulPart} CORRIGÉ --- FICHE D'EXERCICES}\\[3pt]
      {\normalsize\bfseries #1}\\[2pt]
      {\normalsize #2}
    \end{tcolorbox}
  \end{center}}

\newcounter{partie}
\newcommand{\partie}[1]{%
  \stepcounter{partie}%
  \besoinplace{8\baselineskip}\vspace{4pt}%
  \begin{tcolorbox}[enhanced,colback=coulPart!8,colframe=coulPart,boxrule=0pt,
      leftrule=3pt,arc=0pt,left=6pt,right=4pt,top=3pt,bottom=3pt,
      before skip=6pt,after skip=2pt]
    \bfseries\color{coulPart}\Roman{partie} --- #1\signet{1}{\Roman{partie} --- #1}
  \end{tcolorbox}}
\newcommand{\partiesansnum}[1]{%
  \besoinplace{8\baselineskip}\vspace{4pt}%
  \begin{tcolorbox}[enhanced,colback=coulPart!8,colframe=coulPart,boxrule=0pt,
      leftrule=3pt,arc=0pt,left=6pt,right=4pt,top=3pt,bottom=3pt,
      before skip=6pt,after skip=2pt]
    \bfseries\color{coulPart}#1\signet{1}{#1}
  \end{tcolorbox}}

\setlist[questions,1]{label=\arabic*\textdegree),leftmargin=*,itemsep=3pt,topsep=3pt}
\setlist[questions,2]{label=\alph*),leftmargin=*,itemsep=2pt,topsep=2pt}
\setlist[questions,3]{label=\roman*),leftmargin=*,itemsep=1pt,topsep=1pt}
\newlist{expr}{enumerate}{1}
\setlist[expr]{label={\Alph*\ $=$},leftmargin=*,itemsep=2pt,topsep=3pt}

\newcommand{\res}[1]{{\color{coulRep}#1}}
\newcommand{\rem}[1]{\par\smallskip{\small\itshape\color{black!60}$\rightarrow$ #1}}
\newenvironment{calculs}{\par\smallskip\noindent$\displaystyle\begin{aligned}[t]}%
  {\end{aligned}$\par\smallskip}
\newcommand{\justif}[1]{\quad\text{\small\itshape\color{black!60}#1}}
\newcommand{\cas}[1]{\par\smallskip\noindent\textbf{\color{coulExo}#1}\nobreak}
\newcommand{\explique}[2]{\par\smallskip\noindent
  \begin{minipage}[t]{0.6\linewidth}\raggedright #1\end{minipage}\hfill
  \begin{minipage}[t]{0.37\linewidth}\raggedright\leavevmode
    {\small\itshape\color{black!60}#2}\end{minipage}\par}

\setlength{\parskip}{2pt}

\makeatother
\entete{Chapitre 1 --- Fractions et puissances --- corrigé}{Seconde}
\setlength{\parindent}{0pt}
\setlength{\jot}{5pt}

\begin{document}

\entetefiche{Chapitre 1 --- Fractions et puissances}{Seconde}

\partie{Nombres relatifs et priorités}

\begin{exo}[Calculs sur les relatifs]
\begin{calculs}
A &= -9-7 &&\justif{deux nombres négatifs : on ajoute leurs distances à zéro, $9$ et $7$}\\
A &= \res{-16} &&\justif{on garde le signe moins}
\end{calculs}
\begin{calculs}
B &= 14-(-11) &&\justif{on recopie}\\
B &= 14+11 &&\justif{soustraire $-11$, c'est ajouter son opposé $11$}\\
B &= \res{25} &&
\end{calculs}
\begin{calculs}
C &= -6\times 7 &&\justif{signes contraires : le produit est négatif}\\
C &= \res{-42} &&\justif{$6\times 7=42$}
\end{calculs}
\begin{calculs}
D &= (-8)\times(-5) &&\justif{même signe : le produit est positif}\\
D &= \res{40} &&\justif{$8\times 5=40$}
\end{calculs}
\begin{calculs}
E &= \dfrac{-(-27)}{-9} &&\justif{on recopie}\\
E &= \dfrac{27}{-9} &&\justif{l'opposé de $-27$ est $27$}\\
E &= \res{-3} &&\justif{signes contraires : le quotient est négatif ; $27\div 9=3$}
\end{calculs}
\begin{calculs}
F &= -13+(-4) &&\justif{ajouter $-4$, c'est soustraire $4$}\\
F &= -13-4 &&\\
F &= \res{-17} &&\justif{deux nombres négatifs : on ajoute $13$ et $4$}
\end{calculs}
\rem{Erreur fréquente : écrire $-9-7=16$ en pensant \og moins par moins donne plus\fg.
Cette règle vaut pour un \emph{produit} ou un \emph{quotient}, jamais pour une somme :
dans $A$, on perd $9$ puis encore $7$.}
\end{exo}

\begin{exo}[Priorités de calcul]
\begin{calculs}
A &= 8^{2}-3\times(-6) &&\justif{on recopie}\\
A &= 64-3\times(-6) &&\justif{la puissance d'abord : $8^{2}=8\times 8=64$}\\
A &= 64-(-18) &&\justif{puis le produit : signes contraires, négatif}\\
A &= 64+18 &&\justif{soustraire $-18$, c'est ajouter $18$}\\
A &= \res{82} &&
\end{calculs}
\begin{calculs}
B &= -7-4^{2} &&\justif{la puissance passe avant la soustraction}\\
B &= -7-16 &&\justif{$4^{2}=16$ ; le moins devant reste une soustraction}\\
B &= \res{-23} &&
\end{calculs}
\rem{Dans $B$, le carré ne porte que sur $4$ : on ne calcule pas $(-4)^{2}$.}
\begin{calculs}
C &= 6+(-9+5)\times 4 &&\justif{la parenthèse d'abord}\\
C &= 6+(-4)\times 4 &&\justif{$-9+5=-4$}\\
C &= 6+(-16) &&\justif{puis le produit, avant l'addition}\\
C &= \res{-10} &&\justif{signes contraires : $16-6=10$, signe de $-16$}
\end{calculs}
\rem{Erreur fréquente : calculer $6+(-4)=2$, puis $2\times 4=8$. Le produit passe
avant l'addition.}
\begin{calculs}
D &= -36\div(-4)\times 2^{2} &&\justif{la puissance d'abord}\\
D &= -36\div(-4)\times 4 &&\justif{division et produit : de gauche à droite}\\
D &= 9\times 4 &&\justif{$(-36)\div(-4)=9$ : même signe, positif}\\
D &= \res{36} &&
\end{calculs}
\rem{Dans $D$, on ne commence pas par $(-4)\times 4$ : division et multiplication ont
la même priorité, on les fait dans l'ordre où elles sont écrites.}
\end{exo}

\begin{exo}[Parenthèses]
\begin{calculs}
A &= 20-4\times(3-8) &&\justif{la parenthèse d'abord}\\
A &= 20-4\times(-5) &&\justif{$3-8=-5$}\\
A &= 20-(-20) &&\justif{puis le produit : $4\times(-5)=-20$}\\
A &= 20+20 &&\justif{soustraire $-20$, c'est ajouter $20$}\\
A &= \res{40} &&
\end{calculs}
\rem{Erreur fréquente : calculer $20-4=16$ avant le produit. Le $4$ est d'abord
multiplié par la parenthèse.}
\begin{calculs}
B &= 5\times(-3)+2^{3} &&\justif{on recopie}\\
B &= 5\times(-3)+8 &&\justif{la puissance : $2^{3}=2\times 2\times 2=8$, et non $2\times 3$}\\
B &= -15+8 &&\justif{puis le produit}\\
B &= \res{-7} &&\justif{signes contraires : $15-8=7$, signe de $-15$}
\end{calculs}
\begin{calculs}
C &= -48\div\bigl(2\times(-3)\bigr) &&\justif{la parenthèse d'abord}\\
C &= -48\div(-6) &&\justif{$2\times(-3)=-6$}\\
C &= \res{8} &&\justif{même signe : le quotient est positif}
\end{calculs}
\begin{calculs}
D &= 3-2\times\bigl(5-(1-6)\bigr) &&\justif{la parenthèse la plus intérieure d'abord}\\
D &= 3-2\times\bigl(5-(-5)\bigr) &&\justif{$1-6=-5$}\\
D &= 3-2\times 10 &&\justif{$5-(-5)=5+5=10$}\\
D &= 3-20 &&\justif{le produit avant la soustraction}\\
D &= \res{-17} &&
\end{calculs}
\end{exo}

\partie{Additionner et soustraire des fractions}

\begin{exo}[Même dénominateur]
\begin{calculs}
A &= \dfrac{4}{5}+\dfrac{3}{5} &&\justif{même dénominateur : on additionne les numérateurs}\\
A &= \dfrac{4+3}{5} &&\justif{le dénominateur ne change pas}\\
A &= \res{\dfrac{7}{5}} &&\justif{$7$ et $5$ n'ont pas de diviseur commun}
\end{calculs}
\begin{calculs}
B &= \dfrac{-2}{9}+\dfrac{7}{9} &&\\
B &= \dfrac{-2+7}{9} &&\\
B &= \res{\dfrac{5}{9}} &&\justif{$-2+7=5$}
\end{calculs}
\begin{calculs}
C &= \dfrac{5}{7}-\dfrac{9}{7} &&\\
C &= \dfrac{5-9}{7} &&\justif{on soustrait les numérateurs}\\
C &= \res{-\dfrac{4}{7}} &&\justif{$5-9=-4$ ; le signe moins s'écrit devant la fraction}
\end{calculs}
\begin{calculs}
D &= \dfrac{-3}{11}-\dfrac{-8}{11} &&\\
D &= \dfrac{-3-(-8)}{11} &&\justif{parenthèses autour de $-8$}\\
D &= \dfrac{-3+8}{11} &&\justif{soustraire $-8$, c'est ajouter $8$}\\
D &= \res{\dfrac{5}{11}} &&
\end{calculs}
\begin{calculs}
E &= \dfrac{7}{12}+\dfrac{-1}{12} &&\\
E &= \dfrac{7-1}{12} &&\justif{ajouter $-1$, c'est soustraire $1$}\\
E &= \dfrac{6}{12} &&\justif{$6$ et $12$ sont divisibles par $6$ : on simplifie}\\
E &= \dfrac{1\times\barre{6}}{2\times\barre{6}} &&\justif{$6=1\times 6$ et $12=2\times 6$}\\
E &= \res{\dfrac{1}{2}} &&
\end{calculs}
\begin{calculs}
F &= \dfrac{-7}{4}-\dfrac{5}{4} &&\\
F &= \dfrac{-7-5}{4} &&\\
F &= \dfrac{-12}{4} &&\justif{deux nombres négatifs : $-7-5=-12$}\\
F &= \res{-3} &&\justif{$12\div 4=3$ : le résultat est un entier}
\end{calculs}
\rem{On n'additionne jamais les dénominateurs : $\frac45+\frac35$ ne vaut pas
$\frac{7}{10}$. Des cinquièmes plus des cinquièmes donnent des cinquièmes. On pense
ensuite à simplifier ($E$ et $F$).}
\end{exo}

\begin{exo}[Dénominateurs différents]
\explique{Pour chaque calcul, on cherche d'abord un dénominateur commun : le plus petit
multiple commun aux deux dénominateurs.}{on récite la table du plus grand dénominateur
jusqu'à tomber sur un multiple de l'autre}
\begin{calculs}
A &= \dfrac{2}{5}+\dfrac{3}{4} &&\justif{$5$ et $4$ : dénominateur commun $5\times 4=20$}\\
A &= \dfrac{2\fois{4}}{5\fois{4}}+\dfrac{3\fois{5}}{4\fois{5}} &&\justif{on multiplie en haut et en bas par le même nombre}\\
A &= \dfrac{8}{20}+\dfrac{15}{20} &&\\
A &= \res{\dfrac{23}{20}} &&\justif{$23$ et $20$ n'ont pas de diviseur commun}
\end{calculs}
\begin{calculs}
B &= \dfrac{5}{6}+\dfrac{1}{4} &&\justif{table de $6$ : $6$ puis $12$, qui est dans la table de $4$}\\
B &= \dfrac{5\fois{2}}{6\fois{2}}+\dfrac{1\fois{3}}{4\fois{3}} &&\justif{$6\times 2=12$ et $4\times 3=12$}\\
B &= \dfrac{10}{12}+\dfrac{3}{12} &&\\
B &= \res{\dfrac{13}{12}} &&
\end{calculs}
\rem{Avec le dénominateur $6\times 4=24$, on trouve $\frac{26}{24}$ : c'est juste aussi,
mais il faut ensuite simplifier par $2$. Le plus petit dénominateur commun évite cette
étape.}
\begin{calculs}
C &= \dfrac{3}{4}+\dfrac{1}{6} &&\justif{dénominateur commun $12$, comme pour $B$}\\
C &= \dfrac{3\fois{3}}{4\fois{3}}+\dfrac{1\fois{2}}{6\fois{2}} &&\\
C &= \dfrac{9}{12}+\dfrac{2}{12} &&\\
C &= \res{\dfrac{11}{12}} &&
\end{calculs}
\begin{calculs}
D &= \dfrac{3}{8}-\dfrac{5}{6} &&\justif{table de $8$ : $8$ ; $16$ ; $24$, qui est dans la table de $6$}\\
D &= \dfrac{3\fois{3}}{8\fois{3}}-\dfrac{5\fois{4}}{6\fois{4}} &&\\
D &= \dfrac{9}{24}-\dfrac{20}{24} &&\\
D &= \dfrac{9-20}{24} &&\\
D &= \res{-\dfrac{11}{24}} &&\justif{$9-20=-11$}
\end{calculs}
\begin{calculs}
E &= -\dfrac{2}{3}-\dfrac{4}{5} &&\justif{dénominateur commun $3\times 5=15$}\\
E &= -\dfrac{2\fois{5}}{3\fois{5}}-\dfrac{4\fois{3}}{5\fois{3}} &&\\
E &= -\dfrac{10}{15}-\dfrac{12}{15} &&\\
E &= \dfrac{-10-12}{15} &&\justif{deux nombres négatifs : on ajoute $10$ et $12$}\\
E &= \res{-\dfrac{22}{15}} &&
\end{calculs}
\begin{calculs}
F &= -\dfrac{7}{6}+\dfrac{5}{4} &&\justif{dénominateur commun $12$}\\
F &= -\dfrac{7\fois{2}}{6\fois{2}}+\dfrac{5\fois{3}}{4\fois{3}} &&\\
F &= -\dfrac{14}{12}+\dfrac{15}{12} &&\\
F &= \dfrac{-14+15}{12} &&\\
F &= \res{\dfrac{1}{12}} &&\justif{$-14+15=1$}
\end{calculs}
\end{exo}

\begin{exo}[Avec un nombre entier]
\explique{Un entier est une fraction de dénominateur $1$ : $3=\dfrac{3}{1}$. On l'écrit
avec le dénominateur de l'autre fraction.}{même méthode que pour deux fractions}
\begin{calculs}
A &= 3-\dfrac{5}{4} &&\justif{on recopie}\\
A &= \dfrac{3\fois{4}}{1\fois{4}}-\dfrac{5}{4} &&\justif{$3=\frac31$, écrit avec le dénominateur $4$}\\
A &= \dfrac{12}{4}-\dfrac{5}{4} &&\\
A &= \res{\dfrac{7}{4}} &&\justif{$12-5=7$}
\end{calculs}
\begin{calculs}
B &= -2-\dfrac{3}{5} &&\\
B &= -\dfrac{2\fois{5}}{1\fois{5}}-\dfrac{3}{5} &&\\
B &= -\dfrac{10}{5}-\dfrac{3}{5} &&\\
B &= \dfrac{-10-3}{5} &&\justif{deux nombres négatifs : on ajoute $10$ et $3$}\\
B &= \res{-\dfrac{13}{5}} &&
\end{calculs}
\begin{calculs}
C &= -4+\dfrac{17}{6} &&\\
C &= -\dfrac{4\fois{6}}{1\fois{6}}+\dfrac{17}{6} &&\\
C &= -\dfrac{24}{6}+\dfrac{17}{6} &&\\
C &= \dfrac{-24+17}{6} &&\justif{signes contraires : $24-17=7$, signe de $-24$}\\
C &= \res{-\dfrac{7}{6}} &&
\end{calculs}
\begin{calculs}
D &= \dfrac{9}{2}-4 &&\\
D &= \dfrac{9}{2}-\dfrac{4\fois{2}}{1\fois{2}} &&\\
D &= \dfrac{9}{2}-\dfrac{8}{2} &&\\
D &= \res{\dfrac{1}{2}} &&\justif{contrôle : $4{,}5-4=0{,}5$}
\end{calculs}
\rem{Erreur fréquente : écrire $3-\frac54=\frac{3-5}{4}$. On ne retire pas un entier au
seul numérateur : il faut d'abord l'écrire en quarts.}
\end{exo}

\begin{exo}[Avec une lettre]
\explique{On calcule avec $x$ comme avec un nombre : $3x=\dfrac{3x}{1}$ s'écrit avec le
dénominateur de l'autre fraction.}{même méthode qu'avec un entier}
\begin{calculs}
A &= 3x-\dfrac{x}{4} &&\justif{on recopie}\\
A &= \dfrac{3x\fois{4}}{1\fois{4}}-\dfrac{x}{4} &&\justif{$3x$ écrit avec le dénominateur $4$}\\
A &= \dfrac{12x}{4}-\dfrac{x}{4} &&\\
A &= \dfrac{12x-x}{4} &&\\
A &= \res{\dfrac{11x}{4}} &&\justif{$12$ fois $x$, moins une fois $x$ : $11$ fois $x$}
\end{calculs}
\begin{calculs}
B &= \dfrac{2}{5}x-3 &&\justif{$\frac25x=\frac{2x}{5}$}\\
B &= \dfrac{2x}{5}-\dfrac{3\fois{5}}{1\fois{5}} &&\\
B &= \dfrac{2x}{5}-\dfrac{15}{5} &&\\
B &= \res{\dfrac{2x-15}{5}} &&\justif{$2x$ et $-15$ ne se regroupent pas}
\end{calculs}
\begin{calculs}
C &= 2-\dfrac{x-1}{3} &&\\
C &= \dfrac{2\fois{3}}{1\fois{3}}-\dfrac{x-1}{3} &&\\
C &= \dfrac{6}{3}-\dfrac{x-1}{3} &&\\
C &= \dfrac{6-(x-1)}{3} &&\justif{parenthèses : le moins porte sur $x$ et sur $-1$}\\
C &= \dfrac{6-x+1}{3} &&\justif{$-(x-1)=-x+1$}\\
C &= \res{\dfrac{7-x}{3}} &&
\end{calculs}
\rem{Erreur fréquente : écrire $\frac{6-x-1}{3}$. Le signe moins devant une fraction
s'applique à tout son numérateur (propriété 3 du cours).}
\begin{calculs}
D &= \dfrac{2x+1}{3}-x &&\\
D &= \dfrac{2x+1}{3}-\dfrac{x\fois{3}}{1\fois{3}} &&\\
D &= \dfrac{2x+1}{3}-\dfrac{3x}{3} &&\\
D &= \dfrac{2x+1-3x}{3} &&\justif{le moins ne porte que sur $3x$}\\
D &= \res{\dfrac{1-x}{3}} &&\justif{$2x-3x=-x$}
\end{calculs}
\rem{Pour contrôler, on remplace $x$ par un nombre simple dans l'expression de départ et
dans le résultat. Pour $C$ avec $x=4$ : $2-\frac33=1$ et $\frac{7-4}{3}=1$.}
\end{exo}

\partie{Multiplier et diviser des fractions}

\begin{exo}[Simplifier avant de calculer]
\explique{On fixe d'abord le signe en comptant les signes moins, puis on décompose
chaque nombre en produit de petits facteurs pour barrer les facteurs communs.}{on
n'effectue pas les produits : les nombres resteraient grands}
\begin{calculs}
A &= \dfrac{-4\times 9}{6\times 10} &&\justif{un seul signe moins : le résultat est négatif}\\
A &= -\dfrac{2\times 2\times 3\times 3}{2\times 3\times 2\times 5} &&\justif{$4=2\times 2$ ; $9=3\times 3$ ; $6=2\times 3$ ; $10=2\times 5$}\\
A &= -\dfrac{\barre{2}\times\barre{2}\times\barre{3}\times 3}{\barre{2}\times\barre{3}\times\barre{2}\times 5} &&\justif{on barre deux $2$ et un $3$, en haut et en bas}\\
A &= \res{-\dfrac{3}{5}} &&
\end{calculs}
\begin{calculs}
B &= \dfrac{24\times(-35)}{(-21)\times 10} &&\justif{deux signes moins : le résultat est positif}\\
B &= \dfrac{2\times 3\times 4\times 5\times 7}{3\times 7\times 2\times 5} &&\justif{$24=2\times 3\times 4$ ; $35=5\times 7$ ; $21=3\times 7$ ; $10=2\times 5$}\\
B &= \dfrac{\barre{2}\times\barre{3}\times 4\times\barre{5}\times\barre{7}}{\barre{3}\times\barre{7}\times\barre{2}\times\barre{5}} &&\justif{tout le dénominateur se barre : il reste $1$ en bas}\\
B &= \res{4} &&\justif{$\frac41=4$}
\end{calculs}
\begin{calculs}
C &= \dfrac{26\times(-9)}{15\times 39} &&\justif{un seul signe moins : le résultat est négatif}\\
C &= -\dfrac{2\times 13\times 3\times 3}{3\times 5\times 3\times 13} &&\justif{$26=2\times 13$ et $39=3\times 13$ : $13$ est le facteur commun caché}\\
C &= -\dfrac{2\times\barre{13}\times\barre{3}\times\barre{3}}{\barre{3}\times 5\times\barre{3}\times\barre{13}} &&\\
C &= \res{-\dfrac{2}{5}} &&
\end{calculs}
\begin{calculs}
D &= \dfrac{14\times(-15)}{(-10)\times 21} &&\justif{deux signes moins : le résultat est positif}\\
D &= \dfrac{1\times 2\times 7\times 3\times 5}{2\times 5\times 3\times 7} &&\justif{$14=2\times 7$ ; $15=3\times 5$ ; $10=2\times 5$ ; $21=3\times 7$ ; tout le numérateur va se barrer : on écrit le facteur $1$}\\
D &= \dfrac{1\times\barre{2}\times\barre{7}\times\barre{3}\times\barre{5}}{\barre{2}\times\barre{5}\times\barre{3}\times\barre{7}} &&\justif{tout se simplifie}\\
D &= \res{1} &&\justif{il reste $1$, et non $0$}
\end{calculs}
\rem{En effectuant les produits, $B$ devient $\frac{840}{210}$ : la simplification est
bien plus difficile à voir. On décompose d'abord, on barre ensuite.}
\end{exo}

\begin{exo}[Simplifier avec des lettres]
\begin{calculs}
A &= \dfrac{15\times 8\times 14}{35\times 6\times 4} &&\justif{on recopie}\\
A &= \dfrac{3\times 5\times 2\times 4\times 2\times 7}{5\times 7\times 2\times 3\times 4} &&\justif{$15=3\times 5$ ; $8=2\times 4$ ; $14=2\times 7$ ; $35=5\times 7$ ; $6=2\times 3$}\\
A &= \dfrac{\barre{3}\times\barre{5}\times\barre{2}\times\barre{4}\times 2\times\barre{7}}{\barre{5}\times\barre{7}\times\barre{2}\times\barre{3}\times\barre{4}} &&\justif{il ne reste qu'un $2$ en haut}\\
A &= \res{2} &&
\end{calculs}
\begin{calculs}
B &= \dfrac{3x\times 20}{4\times 5} &&\justif{$3x$ veut dire $3\times x$}\\
B &= \dfrac{3\times x\times 4\times 5}{4\times 5} &&\justif{$20=4\times 5$}\\
B &= \dfrac{3\times x\times\barre{4}\times\barre{5}}{\barre{4}\times\barre{5}} &&\\
B &= \res{3x} &&
\end{calculs}
\begin{calculs}
C &= \dfrac{6a\times 5b}{10b\times 9} &&\\
C &= \dfrac{2\times 3\times a\times 5\times b}{2\times 5\times b\times 3\times 3} &&\justif{une lettre est un facteur comme un nombre}\\
C &= \dfrac{\barre{2}\times\barre{3}\times a\times\barre{5}\times\barre{b}}{\barre{2}\times\barre{5}\times\barre{b}\times\barre{3}\times 3} &&\justif{on barre $b$ : il n'est pas nul}\\
C &= \res{\dfrac{a}{3}} &&
\end{calculs}
\begin{calculs}
D &= \dfrac{4a\times 9}{6\times 2a} &&\\
D &= \dfrac{2\times 2\times a\times 3\times 3}{2\times 3\times 2\times a} &&\justif{$4=2\times 2$ ; $9=3\times 3$ ; $6=2\times 3$}\\
D &= \dfrac{\barre{2}\times\barre{2}\times\barre{a}\times 3\times\barre{3}}{\barre{2}\times\barre{3}\times\barre{2}\times\barre{a}} &&\justif{il ne reste qu'un $3$ en haut}\\
D &= \res{3} &&
\end{calculs}
\rem{On ne barre que des facteurs d'un produit. L'énoncé précise que $a$ et $b$ ne sont
pas nuls : on ne peut pas simplifier par $0$.}
\end{exo}

\begin{exo}[Produits]
\explique{Pour multiplier deux fractions, on fixe le signe, puis on multiplie les
numérateurs entre eux et les dénominateurs entre eux.}{pas besoin de dénominateur commun
pour un produit}
\begin{calculs}
A &= \dfrac{5}{4}\times\dfrac{-3}{7} &&\justif{un seul signe moins : le résultat est négatif}\\
A &= -\dfrac{5\times 3}{4\times 7} &&\justif{numérateurs entre eux, dénominateurs entre eux}\\
A &= \res{-\dfrac{15}{28}} &&\justif{$15=3\times 5$ et $28=4\times 7$ : rien à simplifier}
\end{calculs}
\begin{calculs}
B &= \dfrac{-2}{5}\times\dfrac{4}{3} &&\justif{un seul signe moins}\\
B &= -\dfrac{2\times 4}{5\times 3} &&\\
B &= \res{-\dfrac{8}{15}} &&
\end{calculs}
\begin{calculs}
C &= -5\times\dfrac{3}{8} &&\\
C &= -\dfrac{5}{1}\times\dfrac{3}{8} &&\justif{un entier est une fraction de dénominateur $1$}\\
C &= -\dfrac{5\times 3}{1\times 8} &&\\
C &= \res{-\dfrac{15}{8}} &&
\end{calculs}
\rem{Erreur fréquente : multiplier aussi le dénominateur par $5$. On ne multiplie que le
numérateur, puisque $5=\frac51$.}
\begin{calculs}
D &= \dfrac{-7}{-3}\times\dfrac{2}{-5} &&\justif{trois signes moins : le résultat est négatif}\\
D &= -\dfrac{7\times 2}{3\times 5} &&\\
D &= \res{-\dfrac{14}{15}} &&
\end{calculs}
\begin{calculs}
E &= \dfrac{-4}{5}\times\dfrac{-4}{9} &&\justif{deux signes moins : le résultat est positif}\\
E &= \dfrac{4\times 4}{5\times 9} &&\\
E &= \res{\dfrac{16}{45}} &&
\end{calculs}
\begin{calculs}
F &= \dfrac{3}{-8}\times\dfrac{-5}{-7} &&\justif{trois signes moins : le résultat est négatif}\\
F &= -\dfrac{3\times 5}{8\times 7} &&\\
F &= \res{-\dfrac{15}{56}} &&
\end{calculs}
\rem{Le signe d'un produit se lit en comptant les signes moins : un nombre pair donne un
résultat positif, un nombre impair un résultat négatif.}
\end{exo}

\begin{exo}[Quotients]
\explique{Diviser par une fraction, c'est multiplier par son inverse : on retourne la
fraction placée \emph{après} le signe $\div$.}{$\dfrac{a}{b}\div\dfrac{c}{d}=\dfrac{a}{b}\times\dfrac{d}{c}$}
\begin{calculs}
A &= \dfrac{2}{5}\div\dfrac{3}{7} &&\justif{on recopie}\\
A &= \dfrac{2}{5}\times\dfrac{7}{3} &&\justif{on multiplie par l'inverse de $\frac37$}\\
A &= \dfrac{2\times 7}{5\times 3} &&\\
A &= \res{\dfrac{14}{15}} &&
\end{calculs}
\begin{calculs}
B &= \dfrac{5}{4}\div\dfrac{7}{3} &&\\
B &= \dfrac{5}{4}\times\dfrac{3}{7} &&\justif{l'inverse de $\frac73$ est $\frac37$}\\
B &= \dfrac{5\times 3}{4\times 7} &&\\
B &= \res{\dfrac{15}{28}} &&
\end{calculs}
\begin{calculs}
C &= \dfrac{3}{4}\div 6 &&\\
C &= \dfrac{3}{4}\times\dfrac{1}{6} &&\justif{l'inverse de $6$ est $\frac16$}\\
C &= \dfrac{1\times 3}{4\times 2\times 3} &&\justif{$6=2\times 3$ ; le numérateur va se barrer : facteur $1$}\\
C &= \dfrac{1\times\barre{3}}{4\times 2\times\barre{3}} &&\\
C &= \res{\dfrac{1}{8}} &&
\end{calculs}
\begin{calculs}
D &= 6\div\dfrac{-2}{5} &&\\
D &= 6\times\dfrac{5}{-2} &&\justif{l'inverse de $\frac{-2}{5}$ est $\frac{5}{-2}$ : le signe reste}\\
D &= -\dfrac{6\times 5}{2} &&\justif{un seul signe moins}\\
D &= -\dfrac{\barre{2}\times 3\times 5}{\barre{2}} &&\justif{$6=2\times 3$}\\
D &= \res{-15} &&
\end{calculs}
\begin{calculs}
E &= \dfrac{4}{9}\div\dfrac{-5}{2} &&\\
E &= \dfrac{4}{9}\times\dfrac{2}{-5} &&\\
E &= -\dfrac{4\times 2}{9\times 5} &&\justif{un seul signe moins}\\
E &= \res{-\dfrac{8}{45}} &&
\end{calculs}
\begin{calculs}
F &= 3\div\dfrac{-4}{7} &&\\
F &= 3\times\dfrac{7}{-4} &&\\
F &= -\dfrac{3\times 7}{4} &&\\
F &= \res{-\dfrac{21}{4}} &&
\end{calculs}
\rem{Erreur fréquente : retourner la première fraction, ou les deux. Seul le diviseur est
remplacé par son inverse ; l'inverse d'un nombre négatif est négatif.}
\end{exo}

\partie{Fractions et priorités de calcul}

\begin{exo}[Respecter les priorités]
\begin{calculs}
A &= \dfrac{5}{12}-\dfrac{-3}{4} &&\justif{on recopie}\\
A &= \dfrac{5}{12}+\dfrac{3}{4} &&\justif{soustraire $-\frac34$, c'est ajouter $\frac34$}\\
A &= \dfrac{5}{12}+\dfrac{3\fois{3}}{4\fois{3}} &&\justif{$12$ est un multiple de $4$}\\
A &= \dfrac{5}{12}+\dfrac{9}{12} &&\\
A &= \dfrac{14}{12} &&\\
A &= \dfrac{7\times\barre{2}}{6\times\barre{2}} &&\justif{on simplifie par $2$}\\
A &= \res{\dfrac{7}{6}} &&
\end{calculs}
\begin{calculs}
B &= \dfrac{7}{10}-\dfrac{-2}{5} &&\\
B &= \dfrac{7}{10}+\dfrac{2}{5} &&\justif{soustraire $-\frac25$, c'est ajouter $\frac25$}\\
B &= \dfrac{7}{10}+\dfrac{2\fois{2}}{5\fois{2}} &&\justif{$10$ est un multiple de $5$}\\
B &= \dfrac{7}{10}+\dfrac{4}{10} &&\\
B &= \res{\dfrac{11}{10}} &&
\end{calculs}
\begin{calculs}
C &= -\dfrac{5}{6}+\dfrac{3}{4}-\dfrac{2}{3} &&\justif{$12$ est un multiple de $6$, de $4$ et de $3$}\\
C &= -\dfrac{5\fois{2}}{6\fois{2}}+\dfrac{3\fois{3}}{4\fois{3}}-\dfrac{2\fois{4}}{3\fois{4}} &&\\
C &= -\dfrac{10}{12}+\dfrac{9}{12}-\dfrac{8}{12} &&\\
C &= \dfrac{-10+9-8}{12} &&\justif{de gauche à droite : $-10+9=-1$, puis $-1-8=-9$}\\
C &= -\dfrac{9}{12} &&\\
C &= -\dfrac{3\times\barre{3}}{4\times\barre{3}} &&\justif{on simplifie par $3$}\\
C &= \res{-\dfrac{3}{4}} &&
\end{calculs}
\begin{calculs}
D &= \dfrac{9}{14}+\dfrac{5}{7}\times\dfrac{7}{10} &&\justif{le produit passe avant l'addition}\\
D &= \dfrac{9}{14}+\dfrac{1\times 5\times 7}{7\times 2\times 5} &&\justif{$10=2\times 5$ ; on n'effectue pas le produit}\\
D &= \dfrac{9}{14}+\dfrac{1\times\barre{5}\times\barre{7}}{\barre{7}\times 2\times\barre{5}} &&\justif{on barre $5$ et $7$ : il reste $\frac12$}\\
D &= \dfrac{9}{14}+\dfrac{1\fois{7}}{2\fois{7}} &&\justif{$14$ est un multiple de $2$}\\
D &= \dfrac{9}{14}+\dfrac{7}{14} &&\\
D &= \dfrac{16}{14} &&\\
D &= \dfrac{8\times\barre{2}}{7\times\barre{2}} &&\\
D &= \res{\dfrac{8}{7}} &&
\end{calculs}
\rem{Erreur fréquente dans $D$ : additionner $\frac{9}{14}+\frac57$ d'abord. On
trouverait $\frac{19}{20}$, ce qui est faux : le produit passe avant l'addition.}
\end{exo}

\begin{exo}[Parenthèses et lettres]
\begin{calculs}
A &= \dfrac{4}{5}\left(\dfrac{1}{2}-\dfrac{1}{4}\right) &&\justif{un nombre collé à une parenthèse : c'est un produit}\\
A &= \dfrac{4}{5}\times\left(\dfrac{1\fois{2}}{2\fois{2}}-\dfrac{1}{4}\right) &&\justif{la parenthèse d'abord}\\
A &= \dfrac{4}{5}\times\left(\dfrac{2}{4}-\dfrac{1}{4}\right) &&\\
A &= \dfrac{4}{5}\times\dfrac{1}{4} &&\\
A &= \dfrac{\barre{4}\times 1}{5\times\barre{4}} &&\justif{on barre le $4$}\\
A &= \res{\dfrac{1}{5}} &&
\end{calculs}
\begin{calculs}
B &= \left(\dfrac{2}{3}-\dfrac{1}{6}\right)\times\left(\dfrac{1}{4}+\dfrac{1}{2}\right) &&\justif{chaque parenthèse à part}\\
B &= \left(\dfrac{2\fois{2}}{3\fois{2}}-\dfrac{1}{6}\right)\times\left(\dfrac{1}{4}+\dfrac{1\fois{2}}{2\fois{2}}\right) &&\\
B &= \left(\dfrac{4}{6}-\dfrac{1}{6}\right)\times\left(\dfrac{1}{4}+\dfrac{2}{4}\right) &&\\
B &= \dfrac{3}{6}\times\dfrac{3}{4} &&\\
B &= \dfrac{\barre{3}\times 3}{2\times\barre{3}\times 4} &&\justif{$6=2\times 3$ ; on barre un $3$}\\
B &= \res{\dfrac{3}{8}} &&
\end{calculs}
\begin{calculs}
C &= \dfrac{3x}{4}-\dfrac{x-2}{6} &&\justif{$12$ est un multiple de $4$ et de $6$}\\
C &= \dfrac{3x\fois{3}}{4\fois{3}}-\dfrac{(x-2)\fois{2}}{6\fois{2}} &&\justif{tout le numérateur $x-2$ est multiplié : parenthèses}\\
C &= \dfrac{9x}{12}-\dfrac{2x-4}{12} &&\justif{$(x-2)\times 2=2x-4$}\\
C &= \dfrac{9x-(2x-4)}{12} &&\justif{le moins porte sur tout le numérateur}\\
C &= \dfrac{9x-2x+4}{12} &&\justif{$-(2x-4)=-2x+4$}\\
C &= \res{\dfrac{7x+4}{12}} &&
\end{calculs}
\rem{Erreur fréquente : écrire $9x-2x-4$. Contrôle avec $x=2$ : $\frac{6}{4}-\frac{0}{6}=\frac32$
et $\frac{14+4}{12}=\frac{18}{12}=\frac32$.}
\begin{calculs}
D &= \dfrac{3x}{4}\times\dfrac{8}{9}x &&\justif{un produit : pas de dénominateur commun}\\
D &= \dfrac{3x}{4}\times\dfrac{8x}{9} &&\justif{$\frac89x=\frac{8x}{9}$}\\
D &= \dfrac{3\times x\times 2\times 4\times x}{4\times 3\times 3} &&\justif{$8=2\times 4$ et $9=3\times 3$}\\
D &= \dfrac{\barre{3}\times x\times 2\times\barre{4}\times x}{\barre{4}\times\barre{3}\times 3} &&\\
D &= \res{\dfrac{2x^{2}}{3}} &&\justif{$x\times x=x^{2}$}
\end{calculs}
\end{exo}

\begin{exo}[Fractions de fractions]
\explique{Le grand trait de fraction est une division : $\dfrac{\;\frac{a}{b}\;}{\;\frac{c}{d}\;}=\dfrac{a}{b}\div\dfrac{c}{d}$.
On multiplie donc la fraction du haut par l'inverse de celle du bas.}{propriété 6 du
cours}
\begin{calculs}
A &= \dfrac{\;\dfrac{3}{4}\;}{\;\dfrac{5}{6}\;} &&\justif{on recopie}\\
A &= \dfrac{3}{4}\times\dfrac{6}{5} &&\justif{on multiplie par l'inverse de $\frac56$}\\
A &= \dfrac{3\times 2\times 3}{2\times 2\times 5} &&\justif{$6=2\times 3$ et $4=2\times 2$}\\
A &= \dfrac{3\times\barre{2}\times 3}{\barre{2}\times 2\times 5} &&\\
A &= \res{\dfrac{9}{10}} &&
\end{calculs}
\begin{calculs}
B &= \dfrac{\;\dfrac{6}{-35}\;}{\;\dfrac{-9}{14}\;} &&\justif{deux signes moins : le résultat est positif}\\
B &= \dfrac{6}{35}\times\dfrac{14}{9} &&\justif{on multiplie par l'inverse de $\frac{9}{14}$}\\
B &= \dfrac{2\times 3\times 2\times 7}{5\times 7\times 3\times 3} &&\justif{$6=2\times 3$ ; $14=2\times 7$ ; $35=5\times 7$ ; $9=3\times 3$}\\
B &= \dfrac{2\times\barre{3}\times 2\times\barre{7}}{5\times\barre{7}\times\barre{3}\times 3} &&\\
B &= \res{\dfrac{4}{15}} &&
\end{calculs}
\begin{calculs}
C &= \dfrac{\;\dfrac{7}{4}\;}{-14} &&\justif{un seul signe moins : le résultat est négatif}\\
C &= -\dfrac{7}{4}\times\dfrac{1}{14} &&\justif{diviser par $14$, c'est multiplier par $\frac{1}{14}$}\\
C &= -\dfrac{\barre{7}\times 1}{4\times 2\times\barre{7}} &&\justif{$14=2\times 7$}\\
C &= \res{-\dfrac{1}{8}} &&
\end{calculs}
\explique{Pour $D$, le numérateur et le dénominateur sont des calculs : on les calcule à
part, puis on divise.}{le grand trait joue le rôle de parenthèses}
\begin{calculs}
2-\dfrac{3}{5} &= \dfrac{2\fois{5}}{1\fois{5}}-\dfrac{3}{5} &&\justif{le numérateur}\\
2-\dfrac{3}{5} &= \dfrac{10}{5}-\dfrac{3}{5} &&\\
2-\dfrac{3}{5} &= \dfrac{7}{5} &&
\end{calculs}
\begin{calculs}
\dfrac{3}{4}-\dfrac{2}{5} &= \dfrac{3\fois{5}}{4\fois{5}}-\dfrac{2\fois{4}}{5\fois{4}} &&\justif{le dénominateur ; dénominateur commun $20$}\\
\dfrac{3}{4}-\dfrac{2}{5} &= \dfrac{15}{20}-\dfrac{8}{20} &&\\
\dfrac{3}{4}-\dfrac{2}{5} &= \dfrac{7}{20} &&
\end{calculs}
\begin{calculs}
D &= \dfrac{7}{5}\div\dfrac{7}{20} &&\justif{on divise les deux résultats}\\
D &= \dfrac{7}{5}\times\dfrac{20}{7} &&\\
D &= \dfrac{\barre{7}\times 4\times\barre{5}}{\barre{5}\times\barre{7}} &&\justif{$20=4\times 5$}\\
D &= \res{4} &&
\end{calculs}
\rem{Dans $D$, on ne barre pas les $5$ ni les $3$ des petites fractions : ce sont des
sommes, pas des produits.}
\end{exo}

\partie{Produit en croix}

\begin{exo}[Justifier le produit en croix]
\cas{1)}
\explique{Les fractions $\dfrac{3x-1}{4}$ et $\dfrac{x+3}{2}$ sont égales et leurs
dénominateurs $4$ et $2$ ne sont pas nuls. D'après le produit en croix,
$(3x-1)\times 2=4\times(x+3)$, c'est-à-dire $2\times(3x-1)=4\times(x+3)$.}{si
$\frac{a}{b}=\frac{c}{d}$ avec $b\ne 0$ et $d\ne 0$, alors $a\times d=b\times c$}
\explique{On peut aussi le voir directement : on multiplie les deux membres par
$4\times 2=8$. À gauche, $8\times\dfrac{3x-1}{4}=2\times(3x-1)$ ; à droite,
$8\times\dfrac{x+3}{2}=4\times(x+3)$.}{c'est pour cela que le produit en croix
\og marche\fg}
\cas{2)}
\begin{calculs}
2\times(3x-1) &= 4\times(x+3) &&\justif{on part de l'égalité du 1)}\\
6x-2 &= 4x+12 &&\justif{on développe les deux membres}\\
6x-4x &= 12+2 &&\justif{les $x$ à gauche, les nombres à droite}\\
2x &= 14 &&\\
x &= 7 &&\justif{on divise par $2$}
\end{calculs}
\explique{Vérification : $\dfrac{3\times 7-1}{4}=\dfrac{20}{4}=5$ et
$\dfrac{7+3}{2}=\dfrac{10}{2}=5$.}{on remplace $x$ par $7$ dans chaque membre}
\explique{Les deux membres sont égaux : le nombre $x$ vaut \res{$7$}.}{}
\end{exo}

\begin{exo}[Premières équations]
\explique{Chaque équation est une égalité de deux fractions : le produit en croix fait
disparaître les dénominateurs.}{parenthèses autour de chaque numérateur qui est une
somme}
\cas{a)}
\begin{calculs}
\dfrac{x}{4} &= \dfrac{5}{3} &&\justif{on recopie l'équation}\\
3\times x &= 4\times 5 &&\justif{produit en croix}\\
3x &= 20 &&\\
x &= \dfrac{20}{3} &&\justif{on divise par $3$ ; la fraction reste telle quelle}
\end{calculs}
\explique{La solution est \res{$\dfrac{20}{3}$}.}{on ne l'arrondit pas : $20\div 3$ ne
tombe pas juste}
\cas{b)}
\begin{calculs}
\dfrac{4-x}{3} &= \dfrac{x+2}{5} &&\\
5(4-x) &= 3(x+2) &&\justif{produit en croix}\\
20-5x &= 3x+6 &&\justif{on développe}\\
20-6 &= 3x+5x &&\justif{les $x$ du côté où il y en a le plus}\\
14 &= 8x &&\\
x &= \dfrac{14}{8} &&\justif{on divise par $8$}\\
x &= \dfrac{7}{4} &&\justif{on simplifie par $2$}
\end{calculs}
\explique{La solution est \res{$\dfrac{7}{4}$}.}{contrôle : les deux membres valent
$\frac34$}
\cas{c)}
\begin{calculs}
\dfrac{2+x}{-3} &= \dfrac{5}{6} &&\\
6(2+x) &= -3\times 5 &&\justif{le $-3$ garde son signe}\\
12+6x &= -15 &&\\
6x &= -27 &&\justif{on retire $12$ aux deux membres}\\
x &= -\dfrac{27}{6} &&\\
x &= -\dfrac{9}{2} &&\justif{on simplifie par $3$}
\end{calculs}
\explique{La solution est \res{$-\dfrac{9}{2}$}.}{contrôle : $2-\frac92=-\frac52$ et
$-\frac52\div(-3)=\frac56$}
\cas{d)}
\begin{calculs}
\dfrac{2x-1}{3} &= \dfrac{x+4}{-2} &&\\
-2(2x-1) &= 3(x+4) &&\justif{le $-2$ garde son signe}\\
-4x+2 &= 3x+12 &&\justif{$-2\times(-1)=+2$}\\
2-12 &= 3x+4x &&\\
-10 &= 7x &&\\
x &= -\dfrac{10}{7} &&\justif{on divise par $7$}
\end{calculs}
\explique{La solution est \res{$-\dfrac{10}{7}$}.}{contrôle : les deux membres valent
$-\frac97$}
\rem{Erreur fréquente : oublier les parenthèses, $5\times 4-x$ au lieu de $5(4-x)$. Le
facteur multiplie tout le numérateur.}
\end{exo}

\begin{exo}[S'entraîner]
\cas{a)}
\begin{calculs}
\dfrac{x}{5} &= \dfrac{3}{8} &&\\
8x &= 5\times 3 &&\justif{produit en croix}\\
8x &= 15 &&\\
x &= \dfrac{15}{8} &&
\end{calculs}
\explique{Vérification : $\dfrac{15}{8}\div 5=\dfrac{15}{40}=\dfrac{3}{8}$. La solution
est \res{$\dfrac{15}{8}$}.}{diviser par $5$, c'est multiplier par $\frac15$}
\cas{b)}
\begin{calculs}
\dfrac{4}{9} &= \dfrac{x}{6} &&\justif{l'inconnue peut être à droite}\\
4\times 6 &= 9x &&\justif{produit en croix}\\
24 &= 9x &&\\
x &= \dfrac{24}{9} &&\\
x &= \dfrac{8}{3} &&\justif{on simplifie par $3$}
\end{calculs}
\explique{Vérification : $\dfrac{8}{3}\div 6=\dfrac{8}{18}=\dfrac{4}{9}$. La solution
est \res{$\dfrac{8}{3}$}.}{}
\cas{c)}
\begin{calculs}
\dfrac{x+2}{5} &= \dfrac{2x-3}{4} &&\\
4(x+2) &= 5(2x-3) &&\justif{produit en croix}\\
4x+8 &= 10x-15 &&\justif{on développe}\\
8+15 &= 10x-4x &&\justif{les $x$ à droite, où il y en a le plus}\\
23 &= 6x &&\\
x &= \dfrac{23}{6} &&
\end{calculs}
\explique{Vérification : $x+2=\dfrac{35}{6}$, donc $\dfrac{x+2}{5}=\dfrac{7}{6}$ ;
$2x-3=\dfrac{14}{3}$, donc $\dfrac{2x-3}{4}=\dfrac{14}{12}=\dfrac{7}{6}$. La solution
est \res{$\dfrac{23}{6}$}.}{$2x=\frac{23}{3}$ et $\frac{23}{3}-\frac93=\frac{14}{3}$}
\cas{d)}
\begin{calculs}
\dfrac{2x+1}{3} &= \dfrac{3x-2}{4} &&\\
4(2x+1) &= 3(3x-2) &&\\
8x+4 &= 9x-6 &&\\
4+6 &= 9x-8x &&\\
x &= 10 &&
\end{calculs}
\explique{Vérification : $\dfrac{21}{3}=7$ et $\dfrac{28}{4}=7$. La solution est
\res{$10$}.}{}
\end{exo}

\begin{exo}[Plusieurs fractions (1)]
\explique{Un membre contient une \emph{somme} de fractions : le produit en croix ne
s'applique pas. On multiplie chaque terme des deux membres par un dénominateur
commun.}{y compris les termes sans fraction}
\cas{a)}
\begin{calculs}
x+\dfrac{x}{5} &= 18 &&\justif{on recopie ; dénominateur commun $5$}\\
5\times x+5\times\dfrac{x}{5} &= 5\times 18 &&\justif{chaque terme $\times 5$, $x$ compris}\\
5x+x &= 90 &&\justif{$5\times\frac{x}{5}=x$}\\
6x &= 90 &&\\
x &= 15 &&\justif{$90\div 6=15$}
\end{calculs}
\explique{Vérification : $15+\dfrac{15}{5}=15+3=18$. La solution est \res{$15$}.}{}
\cas{b)}
\begin{calculs}
\dfrac{x}{3}-\dfrac{x}{4} &= 5 &&\justif{dénominateur commun $12$}\\
12\times\dfrac{x}{3}-12\times\dfrac{x}{4} &= 12\times 5 &&\justif{chaque terme $\times 12$}\\
4x-3x &= 60 &&\justif{$12\div 3=4$ et $12\div 4=3$}\\
x &= 60 &&
\end{calculs}
\explique{Vérification : $\dfrac{60}{3}-\dfrac{60}{4}=20-15=5$. La solution est
\res{$60$}.}{}
\cas{c)}
\begin{calculs}
\dfrac{t}{3}-\dfrac{t+1}{5} &= 1 &&\justif{dénominateur commun $15$}\\
5t-3(t+1) &= 15 &&\justif{chaque terme $\times 15$ ; parenthèses à cause du moins}\\
5t-3t-3 &= 15 &&\justif{$-3\times(t+1)=-3t-3$}\\
2t &= 18 &&\justif{on ajoute $3$}\\
t &= 9 &&
\end{calculs}
\explique{Vérification : $\dfrac{9}{3}-\dfrac{10}{5}=3-2=1$. La solution est
\res{$9$}.}{}
\cas{d)}
\begin{calculs}
\dfrac{3t-1}{2}+\dfrac{t-2}{3} &= 2t &&\justif{dénominateur commun $6$}\\
3(3t-1)+2(t-2) &= 12t &&\justif{chaque terme $\times 6$ : $6\div 2=3$ et $6\div 3=2$}\\
9t-3+2t-4 &= 12t &&\justif{on développe}\\
11t-7 &= 12t &&\\
-7 &= 12t-11t &&\justif{les $t$ à droite, où il y en a le plus}\\
t &= -7 &&
\end{calculs}
\explique{Vérification : $\dfrac{-22}{2}+\dfrac{-9}{3}=-11-3=-14$ et
$2\times(-7)=-14$. La solution est \res{$-7$}.}{}
\rem{Erreur fréquente : oublier de multiplier le terme sans fraction ($18$ en a), $1$ en
c), $2t$ en d)). L'égalité n'est conservée que si l'on multiplie tout.}
\end{exo}

\begin{exo}[Plusieurs fractions (2)]
\cas{a)}
\begin{calculs}
\dfrac{y}{2}+\dfrac{3y}{4} &= \dfrac{y}{3}-22 &&\justif{$12$ est un multiple de $2$, de $4$ et de $3$}\\
6y+9y &= 4y-264 &&\justif{chaque terme $\times 12$ ; $12\times 22=264$}\\
15y &= 4y-264 &&\\
11y &= -264 &&\justif{on retire $4y$}\\
y &= -24 &&\justif{$264\div 11=24$}
\end{calculs}
\explique{Vérification : $\dfrac{-24}{2}+\dfrac{-72}{4}=-12-18=-30$ et
$\dfrac{-24}{3}-22=-8-22=-30$. La solution est \res{$-24$}.}{}
\cas{b)}
\begin{calculs}
\dfrac{y+7}{3} &= \dfrac{y-1}{2}-\dfrac{y}{4} &&\justif{une somme à droite : pas de produit en croix ; $\times 12$}\\
4(y+7) &= 6(y-1)-3y &&\justif{$12\div 3=4$ ; $12\div 2=6$ ; $12\div 4=3$}\\
4y+28 &= 6y-6-3y &&\\
4y+28 &= 3y-6 &&\\
4y-3y &= -6-28 &&\\
y &= -34 &&
\end{calculs}
\explique{Vérification : $\dfrac{-27}{3}=-9$ et
$\dfrac{-35}{2}-\dfrac{-34}{4}=-17{,}5+8{,}5=-9$. La solution est \res{$-34$}.}{}
\cas{c)}
\begin{calculs}
\dfrac{z-4}{3}-\dfrac{z-6}{12} &= \dfrac{2z-15}{9} &&\justif{table de $12$ : $12$ ; $24$ ; $36$, multiple de $3$ et de $9$}\\
12(z-4)-3(z-6) &= 4(2z-15) &&\justif{chaque terme $\times 36$ : $36\div 3=12$ ; $36\div 12=3$ ; $36\div 9=4$}\\
12z-48-3z+18 &= 8z-60 &&\justif{attention : $-3\times(-6)=+18$}\\
9z-30 &= 8z-60 &&\\
9z-8z &= -60+30 &&\\
z &= -30 &&
\end{calculs}
\explique{Vérification : $\dfrac{-34}{3}-\dfrac{-36}{12}=-\dfrac{34}{3}+3=-\dfrac{25}{3}$
et $\dfrac{-75}{9}=-\dfrac{25}{3}$. La solution est \res{$-30$}.}{$3=\frac93$}
\cas{d)}
\begin{calculs}
\dfrac{3z+2}{5}-\dfrac{2z-1}{2} &= \dfrac{3-6z}{10} &&\justif{dénominateur commun $10$}\\
2(3z+2)-5(2z-1) &= 3-6z &&\justif{chaque terme $\times 10$}\\
6z+4-10z+5 &= 3-6z &&\justif{attention : $-5\times(-1)=+5$}\\
-4z+9 &= 3-6z &&\\
-4z+6z &= 3-9 &&\justif{on ajoute $6z$ et on retire $9$}\\
2z &= -6 &&\\
z &= -3 &&
\end{calculs}
\explique{Vérification : $\dfrac{-7}{5}-\dfrac{-7}{2}=\dfrac{-14+35}{10}=\dfrac{21}{10}$
et $\dfrac{3+18}{10}=\dfrac{21}{10}$. La solution est \res{$-3$}.}{}
\rem{Le piège de cet exercice est le signe moins devant une fraction : après la
multiplication, il porte sur tout le numérateur, d'où les parenthèses.}
\end{exo}

\partie{Fractions et problèmes}

\begin{exo}[Fraction d'une quantité]
\explique{Prendre les $\dfrac{a}{b}$ \emph{de} quelque chose, c'est multiplier par
$\dfrac{a}{b}$.}{en pratique : on divise par $b$, puis on multiplie par $a$}
\cas{a)}
\begin{calculs}
\text{montant} &= \dfrac{3}{4}\times 120 &&\justif{les trois quarts de $120$}\\
\text{montant} &= 3\times 30 &&\justif{un quart de $120$ : $120\div 4=30$}\\
\text{montant} &= 90 &&
\end{calculs}
\explique{Les trois quarts de $120$~€ font \res{$90$~€}.}{}
\cas{b)}
\begin{calculs}
\text{masse} &= \dfrac{2}{5}\times 35 &&\\
\text{masse} &= 2\times 7 &&\justif{un cinquième de $35$ : $35\div 5=7$}\\
\text{masse} &= 14 &&
\end{calculs}
\explique{Les deux cinquièmes de $35$~kg font \res{$14$~kg}.}{}
\cas{c)}
\begin{calculs}
\text{durée} &= \dfrac{5}{6}\times 90 &&\\
\text{durée} &= 5\times 15 &&\justif{un sixième de $90$ : $90\div 6=15$}\\
\text{durée} &= 75 &&
\end{calculs}
\explique{Les cinq sixièmes de $90$~min font \res{$75$~min}.}{soit $1$~h $15$~min}
\cas{d)}
\begin{calculs}
\text{élèves} &= \dfrac{2}{3}\times\dfrac{3}{4}\times 48 &&\justif{\og les $\frac23$ des $\frac34$ de\fg{} : deux multiplications}\\
\text{élèves} &= \dfrac{1\times\barre{2}\times\barre{3}}{\barre{3}\times\barre{2}\times 2}\times 48 &&\justif{$4=2\times 2$}\\
\text{élèves} &= \dfrac{1}{2}\times 48 &&\justif{les $\frac23$ des $\frac34$, c'est la moitié}\\
\text{élèves} &= 24 &&
\end{calculs}
\explique{Cela fait \res{$24$ élèves}.}{contrôle : les $\frac34$ de $48$ font $36$, et
les $\frac23$ de $36$ font $24$}
\end{exo}

\begin{exo}[Le festival]
\begin{calculs}
\text{places en prévente} &= \dfrac{5}{6}\times 1~200 &&\justif{les $\frac56$ des places}\\
\text{places en prévente} &= 5\times 200 &&\justif{$1~200\div 6=200$}\\
\text{places en prévente} &= 1~000 &&
\end{calculs}
\begin{calculs}
\text{places restantes} &= 1~200-1~000 &&\\
\text{places restantes} &= 200 &&
\end{calculs}
\begin{calculs}
\text{places vendues sur place} &= \dfrac{3}{4}\times 200 &&\justif{les $\frac34$ des places \emph{restantes}, pas du total}\\
\text{places vendues sur place} &= 3\times 50 &&\justif{$200\div 4=50$}\\
\text{places vendues sur place} &= 150 &&
\end{calculs}
\begin{calculs}
\text{places invendues} &= 200-150 &&\\
\text{places invendues} &= 50 &&
\end{calculs}
\explique{Il reste \res{$50$ places invendues}.}{contrôle : $1~000+150+50=1~200$}
\rem{Erreur fréquente : prendre les $\frac34$ des $1~200$ places, soit $900$ places ; on
dépasserait le nombre de places disponibles. Une fraction \og du reste\fg{} porte sur ce
qui reste.}
\end{exo}

\begin{exo}[L'anniversaire de Lina]
\cas{a)}
\explique{La question demande une fraction : on raisonne en fractions de la somme
reçue, qui vaut $1$ (la somme entière).}{les euros viennent au b)}
\begin{calculs}
\text{reste après les jeux} &= 1-\dfrac{1}{3} &&\justif{elle dépense le tiers}\\
\text{reste après les jeux} &= \dfrac{3}{3}-\dfrac{1}{3} &&\justif{$1=\frac33$}\\
\text{reste après les jeux} &= \dfrac{2}{3} &&
\end{calculs}
\begin{calculs}
\text{vêtements} &= \dfrac{3}{8}\times\dfrac{2}{3} &&\justif{les $\frac38$ de ce qui reste}\\
\text{vêtements} &= \dfrac{1\times\barre{3}\times\barre{2}}{\barre{2}\times 4\times\barre{3}} &&\justif{$8=2\times 4$}\\
\text{vêtements} &= \dfrac{1}{4} &&\justif{les vêtements font le quart de la somme reçue}
\end{calculs}
\begin{calculs}
\text{reste final} &= \dfrac{2}{3}-\dfrac{1}{4} &&\\
\text{reste final} &= \dfrac{2\fois{4}}{3\fois{4}}-\dfrac{1\fois{3}}{4\fois{3}} &&\justif{dénominateur commun $12$}\\
\text{reste final} &= \dfrac{8}{12}-\dfrac{3}{12} &&\\
\text{reste final} &= \dfrac{5}{12} &&
\end{calculs}
\explique{Il reste à Lina les \res{$\dfrac{5}{12}$} de la somme reçue.}{}
\cas{b)}
\begin{calculs}
\text{jeux vidéo} &= \dfrac{1}{3}\times 120 = 40 &&\justif{le tiers de $120$}\\
\text{vêtements} &= \dfrac{1}{4}\times 120 = 30 &&\justif{le quart de la somme reçue, trouvé au a)}\\
\text{reste} &= \dfrac{5}{12}\times 120 &&\\
\text{reste} &= 5\times 10 = 50 &&\justif{$120\div 12=10$}
\end{calculs}
\explique{Lina a dépensé \res{$40$~€} en jeux vidéo et \res{$30$~€} en vêtements ; il
lui reste \res{$50$~€}.}{contrôle : $40+30+50=120$}
\rem{Autre méthode, en euros : il reste $120-40=80$~€ après les jeux ; les $\frac38$ de
$80$~€ font $30$~€ ; il reste $80-30=50$~€.}
\end{exo}

\begin{exo}[Le budget d'une association]
\cas{a)}
\begin{calculs}
\text{reste après le matériel} &= 1-\dfrac{2}{5} &&\justif{le budget entier vaut $1$}\\
\text{reste après le matériel} &= \dfrac{5}{5}-\dfrac{2}{5} &&\\
\text{reste après le matériel} &= \dfrac{3}{5} &&
\end{calculs}
\begin{calculs}
\text{sorties} &= \dfrac{3}{4}\times\dfrac{3}{5} &&\justif{les $\frac34$ de ce qui reste}\\
\text{sorties} &= \dfrac{9}{20} &&\justif{rien à simplifier}
\end{calculs}
\begin{calculs}
\text{frais} &= \dfrac{3}{5}-\dfrac{9}{20} &&\justif{ce qui reste après les sorties}\\
\text{frais} &= \dfrac{3\fois{4}}{5\fois{4}}-\dfrac{9}{20} &&\justif{$20$ est un multiple de $5$}\\
\text{frais} &= \dfrac{12}{20}-\dfrac{9}{20} &&\\
\text{frais} &= \dfrac{3}{20} &&
\end{calculs}
\explique{Les sorties représentent les \res{$\dfrac{9}{20}$} du budget, les frais de
fonctionnement les \res{$\dfrac{3}{20}$}.}{contrôle : $\frac25=\frac{8}{20}$ et
$\frac{8}{20}+\frac{9}{20}+\frac{3}{20}=1$}
\cas{b)}
\explique{Les frais de fonctionnement, soit les $\dfrac{3}{20}$ du budget, font
$270$~€. On remonte au budget entier en passant par un vingtième.}{le total est
l'inconnue}
\begin{calculs}
\text{un vingtième du budget} &= 270\div 3 &&\justif{trois vingtièmes font $270$~€}\\
\text{un vingtième du budget} &= 90 &&
\end{calculs}
\begin{calculs}
\text{budget} &= 20\times 90 &&\justif{vingt vingtièmes : le budget entier}\\
\text{budget} &= 1~800 &&
\end{calculs}
\explique{Le budget annuel de l'association est de \res{$1~800$~€}.}{contrôle :
matériel $720$~€, sorties $810$~€, frais $270$~€ ; total $1~800$~€}
\rem{Erreur fréquente : calculer $\frac{3}{20}\times 270$. Ici, $270$~€ n'est pas le
total mais une partie : on cherche le nombre dont les $\frac{3}{20}$ font $270$.}
\end{exo}

\partie{Puissances}

\begin{exo}[Calculer des puissances]
\begin{calculs}
A &= 3^{4} &&\justif{$4$ facteurs égaux à $3$}\\
A &= 3\times 3\times 3\times 3 &&\\
A &= \res{81} &&\justif{$9\times 9=81$}
\end{calculs}
\begin{calculs}
B &= (-2)^{4} &&\\
B &= (-2)\times(-2)\times(-2)\times(-2) &&\\
B &= \res{16} &&\justif{exposant pair : le résultat est positif}
\end{calculs}
\begin{calculs}
C &= 4^{-2} &&\justif{$a^{-n}=\frac{1}{a^{n}}$}\\
C &= \dfrac{1}{4^{2}} &&\\
C &= \res{\dfrac{1}{16}} &&
\end{calculs}
\begin{calculs}
D &= (-3)^{-3} &&\\
D &= \dfrac{1}{(-3)^{3}} &&\justif{on garde les parenthèses}\\
D &= \dfrac{1}{-27} &&\justif{exposant impair : $(-3)^{3}$ est négatif}\\
D &= \res{-\dfrac{1}{27}} &&
\end{calculs}
\begin{calculs}
E &= 10^{-3} &&\\
E &= \dfrac{1}{10^{3}} &&\\
E &= \dfrac{1}{1~000} &&\\
E &= \res{0{,}001} &&
\end{calculs}
\begin{calculs}
F &= 25^{0} &&\\
F &= \res{1} &&\justif{$a^{0}=1$ pour tout nombre $a$ non nul}
\end{calculs}
\begin{calculs}
G &= 7^{-1} &&\\
G &= \res{\dfrac{1}{7}} &&\justif{$a^{-1}$ est l'inverse de $a$}
\end{calculs}
\begin{calculs}
H &= (-9)^{1} &&\\
H &= \res{-9} &&\justif{un seul facteur}
\end{calculs}
\rem{Erreur fréquente : croire que $4^{-2}$ est négatif, ou le confondre avec
$4\times(-2)$. Un exposant négatif indique une division : $4^{-2}$ est l'inverse de
$4^{2}$.}
\end{exo}

\begin{exo}[Écriture décimale]
\begin{calculs}
A &= 10^{-3} &&\justif{$3$ chiffres après la virgule, le dernier est $1$}\\
A &= \res{0{,}001} &&
\end{calculs}
\begin{calculs}
B &= 10^{-7} &&\justif{$7$ chiffres après la virgule : six zéros puis $1$}\\
B &= \res{0{,}000~000~1} &&
\end{calculs}
\begin{calculs}
C &= 10^{5} &&\justif{$1$ suivi de $5$ zéros}\\
C &= \res{100~000} &&
\end{calculs}
\begin{calculs}
D &= 10^{0} &&\\
D &= \res{1} &&\justif{$a^{0}=1$}
\end{calculs}
\rem{Erreur fréquente : écrire $7$ zéros après la virgule pour $10^{-7}$. L'exposant
compte les chiffres après la virgule, $1$ compris.}
\end{exo}

\begin{exo}[Puissances et priorités]
\begin{calculs}
A &= 2+5^{2} &&\justif{la puissance avant l'addition}\\
A &= 2+25 &&\\
A &= \res{27} &&
\end{calculs}
\begin{calculs}
B &= (2+5)^{2} &&\justif{la parenthèse d'abord}\\
B &= 7^{2} &&\\
B &= \res{49} &&
\end{calculs}
\rem{$A$ et $B$ ne diffèrent que par les parenthèses, et les résultats changent. De même,
$(2+5)^{2}$ ne vaut pas $2^{2}+5^{2}=29$.}
\begin{calculs}
C &= 2-5^{2} &&\justif{le carré ne porte que sur $5$}\\
C &= 2-25 &&\\
C &= \res{-23} &&
\end{calculs}
\begin{calculs}
D &= 3^{3}+2^{4} &&\\
D &= 27+16 &&\justif{$3\times 3\times 3=27$ et $2\times 2\times 2\times 2=16$}\\
D &= \res{43} &&
\end{calculs}
\begin{calculs}
E &= 2^{4}\times 3^{2} &&\justif{les puissances d'abord}\\
E &= 16\times 9 &&\\
E &= \res{144} &&
\end{calculs}
\begin{calculs}
F &= (-3)^{3}+2^{5} &&\\
F &= -27+32 &&\justif{exposant impair : $(-3)^{3}=-27$ ; $2^{5}=32$}\\
F &= \res{5} &&
\end{calculs}
\end{exo}

\begin{exo}[Puissances et priorités (suite)]
\begin{calculs}
A &= (3^{3}-4^{2})\times 10^{2} &&\justif{dans la parenthèse, les puissances d'abord}\\
A &= (27-16)\times 10^{2} &&\\
A &= 11\times 100 &&\justif{$10^{2}=100$}\\
A &= \res{1~100} &&
\end{calculs}
\begin{calculs}
B &= (-3)^{2}+4^{2}\times 10^{-1} &&\\
B &= 9+16\times 0{,}1 &&\justif{exposant pair : $(-3)^{2}=9$ ; $10^{-1}=0{,}1$}\\
B &= 9+1{,}6 &&\justif{le produit avant l'addition}\\
B &= \res{10{,}6} &&
\end{calculs}
\begin{calculs}
C &= (-4)^{2}-10^{2}\div 5^{2} &&\justif{les puissances d'abord}\\
C &= 16-100\div 25 &&\\
C &= 16-4 &&\justif{la division avant la soustraction}\\
C &= \res{12} &&
\end{calculs}
\begin{calculs}
D &= 10^{2}\times(-1)^{7}-(-2)^{5} &&\\
D &= 100\times(-1)-(-32) &&\justif{exposants impairs : $(-1)^{7}=-1$ et $(-2)^{5}=-32$}\\
D &= -100+32 &&\justif{soustraire $-32$, c'est ajouter $32$}\\
D &= \res{-68} &&
\end{calculs}
\end{exo}

\begin{exo}[Des calculs plus longs]
\begin{calculs}
A &= 20-3^{2}\times(6-2\times 4)^{2} &&\justif{dans la parenthèse, le produit d'abord}\\
A &= 20-3^{2}\times(6-8)^{2} &&\\
A &= 20-3^{2}\times(-2)^{2} &&\\
A &= 20-9\times 4 &&\justif{les puissances : $(-2)^{2}=4$}\\
A &= 20-36 &&\justif{le produit avant la soustraction}\\
A &= \res{-16} &&
\end{calculs}
\begin{calculs}
B &= \dfrac{5^{3}-5\times 4^{2}}{3^{4}-6\times 3^{2}} &&\justif{numérateur et dénominateur à part, comme entre parenthèses}\\
B &= \dfrac{125-5\times 16}{81-6\times 9} &&\justif{les puissances}\\
B &= \dfrac{125-80}{81-54} &&\justif{les produits}\\
B &= \dfrac{45}{27} &&\\
B &= \dfrac{5\times\barre{9}}{3\times\barre{9}} &&\justif{on simplifie par $9$}\\
B &= \res{\dfrac{5}{3}} &&
\end{calculs}
\begin{calculs}
C &= \dfrac{6\times 7-2^{5}}{5\times 2^{3}-(3\times 2)^{2}} &&\\
C &= \dfrac{42-32}{5\times 8-6^{2}} &&\justif{$2^{5}=32$ ; $2^{3}=8$ ; la parenthèse : $3\times 2=6$}\\
C &= \dfrac{10}{40-36} &&\\
C &= \dfrac{10}{4} &&\\
C &= \dfrac{5\times\barre{2}}{2\times\barre{2}} &&\\
C &= \res{\dfrac{5}{2}} &&
\end{calculs}
\rem{Erreur fréquente dans $B$ : barrer les $5$ ou les $3$ de part et d'autre du trait.
Le numérateur et le dénominateur sont des différences : on les calcule avant de
simplifier.}
\end{exo}

\begin{exo}[Des résultats en fraction]
\begin{calculs}
A &= (-3)^{-2}+2\times 10^{-1} &&\\
A &= \dfrac{1}{(-3)^{2}}+2\times\dfrac{1}{10} &&\justif{$a^{-n}=\frac{1}{a^{n}}$}\\
A &= \dfrac{1}{9}+\dfrac{2}{10} &&\\
A &= \dfrac{1}{9}+\dfrac{1}{5} &&\justif{$\frac{2}{10}=\frac15$}\\
A &= \dfrac{1\fois{5}}{9\fois{5}}+\dfrac{1\fois{9}}{5\fois{9}} &&\justif{dénominateur commun $45$}\\
A &= \dfrac{5}{45}+\dfrac{9}{45} &&\\
A &= \res{\dfrac{14}{45}} &&
\end{calculs}
\begin{calculs}
B &= \dfrac{3^{2}-(-2)^{3}}{10^{-1}} &&\\
B &= \dfrac{9-(-8)}{10^{-1}} &&\justif{exposant impair : $(-2)^{3}=-8$}\\
B &= \dfrac{17}{10^{-1}} &&\justif{$9+8=17$}\\
B &= 17\div\dfrac{1}{10} &&\justif{$10^{-1}=\frac{1}{10}$}\\
B &= 17\times 10 &&\justif{diviser par $\frac{1}{10}$, c'est multiplier par $10$}\\
B &= \res{170} &&
\end{calculs}
\begin{calculs}
C &= 2^{-2}-10^{-1} &&\\
C &= \dfrac{1}{4}-\dfrac{1}{10} &&\\
C &= \dfrac{1\fois{5}}{4\fois{5}}-\dfrac{1\fois{2}}{10\fois{2}} &&\justif{dénominateur commun $20$}\\
C &= \dfrac{5}{20}-\dfrac{2}{20} &&\\
C &= \res{\dfrac{3}{20}} &&
\end{calculs}
\begin{calculs}
D &= (2^{5}-4^{2})\times(-4)^{-2} &&\\
D &= (32-16)\times\dfrac{1}{(-4)^{2}} &&\\
D &= 16\times\dfrac{1}{16} &&\justif{exposant pair : $(-4)^{2}=16$}\\
D &= \res{1} &&
\end{calculs}
\end{exo}

\partie{Règles de calcul sur les puissances}

\begin{exo}[Produits]
\explique{Les puissances ont la même base : on ajoute les exposants.}{$a^{n}\times
a^{p}=a^{n+p}$ : on compte $n+p$ facteurs}
\begin{calculs}
A &= 3^{4}\times 3^{7} &&\\
A &= 3^{4+7} &&\\
A &= \res{3^{11}} &&
\end{calculs}
\begin{calculs}
B &= (-5)^{3}\times(-5)^{5} &&\justif{la base est $-5$ : on garde les parenthèses}\\
B &= (-5)^{3+5} &&\\
B &= \res{(-5)^{8}} &&
\end{calculs}
\begin{calculs}
C &= 7^{2}\times 7^{-6} &&\\
C &= 7^{2+(-6)} &&\\
C &= \res{7^{-4}} &&
\end{calculs}
\begin{calculs}
D &= x^{3}\times x\times x^{5} &&\justif{$x=x^{1}$}\\
D &= x^{3+1+5} &&\\
D &= \res{x^{9}} &&
\end{calculs}
\rem{Erreur fréquente : multiplier les exposants ($3^{28}$) ou les bases ($9^{11}$).
$3^{4}\times 3^{7}$ compte $4+7$ facteurs égaux à $3$.}
\end{exo}

\begin{exo}[Quotients]
\explique{Les puissances ont la même base : on soustrait l'exposant du bas à celui du
haut.}{$\dfrac{a^{n}}{a^{p}}=a^{n-p}$}
\begin{calculs}
A &= \dfrac{6^{9}}{6^{4}} &&\\
A &= 6^{9-4} &&\\
A &= \res{6^{5}} &&
\end{calculs}
\begin{calculs}
B &= \dfrac{(-3)^{8}}{(-3)^{5}} &&\\
B &= (-3)^{8-5} &&\\
B &= \res{(-3)^{3}} &&
\end{calculs}
\begin{calculs}
C &= \dfrac{y^{12}}{y^{9}} &&\\
C &= y^{12-9} &&\\
C &= \res{y^{3}} &&
\end{calculs}
\begin{calculs}
D &= \dfrac{2^{-4}}{2^{3}} &&\\
D &= 2^{-4-3} &&\\
D &= \res{2^{-7}} &&
\end{calculs}
\begin{calculs}
E &= \dfrac{4^{-1}}{4^{6}} &&\\
E &= 4^{-1-6} &&\\
E &= \res{4^{-7}} &&
\end{calculs}
\begin{calculs}
F &= \dfrac{x^{-5}}{x^{-9}} &&\\
F &= x^{-5-(-9)} &&\justif{parenthèses autour de l'exposant négatif}\\
F &= x^{-5+9} &&\justif{soustraire $-9$, c'est ajouter $9$}\\
F &= \res{x^{4}} &&
\end{calculs}
\rem{L'exposant du bas est toujours soustrait, même quand il est négatif ($F$).}
\end{exo}

\begin{exo}[Produits et quotients]
\explique{On regroupe d'abord le numérateur et le dénominateur, chacun en une seule
puissance, puis on divise.}{d'abord $a^{n}\times a^{p}$, puis $\frac{a^{n}}{a^{p}}$}
\begin{calculs}
A &= \dfrac{3^{10}\times 3^{4}}{3^{5}\times 3^{6}} &&\\
A &= \dfrac{3^{14}}{3^{11}} &&\justif{$10+4=14$ et $5+6=11$}\\
A &= 3^{14-11} &&\\
A &= \res{3^{3}} &&\justif{soit $27$, non demandé}
\end{calculs}
\begin{calculs}
B &= \dfrac{x^{5}}{x^{7}\times x^{4}} &&\\
B &= \dfrac{x^{5}}{x^{11}} &&\\
B &= x^{5-11} &&\\
B &= \res{x^{-6}} &&
\end{calculs}
\begin{calculs}
C &= \dfrac{2^{4}\times 2^{7}}{2^{3}\times 2^{-2}} &&\\
C &= \dfrac{2^{11}}{2^{1}} &&\justif{$4+7=11$ et $3+(-2)=1$}\\
C &= 2^{11-1} &&\\
C &= \res{2^{10}} &&
\end{calculs}
\begin{calculs}
D &= \dfrac{(y^{2})^{5}}{y^{3}\times y^{4}} &&\\
D &= \dfrac{y^{10}}{y^{7}} &&\justif{$(a^{n})^{p}=a^{n\times p}$ : $2\times 5=10$}\\
D &= y^{10-7} &&\\
D &= \res{y^{3}} &&
\end{calculs}
\end{exo}

\begin{exo}[Toutes les règles]
\explique{On développe d'abord les puissances de puissances et les puissances de
produits, puis on traite chaque base à part.}{même démarche que la méthode 12 du cours}
\begin{calculs}
A &= \dfrac{(5^{2})^{-2}\times 3}{(5\times 3)^{-3}} &&\justif{on recopie}\\
A &= \dfrac{5^{-4}\times 3^{1}}{5^{-3}\times 3^{-3}} &&\justif{$(5^{2})^{-2}=5^{2\times(-2)}$ ; $(5\times 3)^{-3}=5^{-3}\times 3^{-3}$}\\
A &= 5^{-4-(-3)}\times 3^{1-(-3)} &&\justif{chaque base à part ; on soustrait les exposants}\\
A &= 5^{-1}\times 3^{4} &&\justif{$-4+3=-1$ et $1+3=4$}\\
A &= \dfrac{1}{5}\times 81 &&\\
A &= \res{\dfrac{81}{5}} &&
\end{calculs}
\begin{calculs}
B &= \dfrac{3^{4}\times 10^{5}\times 3^{-4}\times 10^{2}}{(3^{3})^{-1}\times(10^{2})^{-4}} &&\\
B &= \dfrac{3^{4}\times 3^{-4}\times 10^{5}\times 10^{2}}{3^{-3}\times 10^{-8}} &&\justif{on regroupe les bases ; $3\times(-1)=-3$ et $2\times(-4)=-8$}\\
B &= \dfrac{3^{0}\times 10^{7}}{3^{-3}\times 10^{-8}} &&\justif{$4+(-4)=0$ et $5+2=7$}\\
B &= 3^{0-(-3)}\times 10^{7-(-8)} &&\\
B &= 3^{3}\times 10^{15} &&\\
B &= \res{27\times 10^{15}} &&
\end{calculs}
\rem{$3^{4}\times 3^{-4}=3^{0}=1$ : les deux facteurs se compensent. Attention aux
exposants négatifs du dénominateur, qu'on soustrait entre parenthèses.}
\end{exo}

\begin{exo}[Calculer et simplifier]
\begin{calculs}
A &= \dfrac{2^{6}\times 10^{-3}\times(2^{3})^{2}\times(10^{-1})^{4}}{2^{10}\times 10^{-5}} &&\\
A &= \dfrac{2^{6}\times 2^{6}\times 10^{-3}\times 10^{-4}}{2^{10}\times 10^{-5}} &&\justif{$(2^{3})^{2}=2^{6}$ et $(10^{-1})^{4}=10^{-4}$}\\
A &= \dfrac{2^{12}\times 10^{-7}}{2^{10}\times 10^{-5}} &&\justif{on ajoute les exposants}\\
A &= 2^{12-10}\times 10^{-7-(-5)} &&\justif{chaque base à part}\\
A &= 2^{2}\times 10^{-2} &&\\
A &= \res{4\times 10^{-2}} &&\justif{soit $0{,}04$}
\end{calculs}
\begin{calculs}
B &= \dfrac{b^{2}\times(b^{-3})^{3}}{(b^{4})^{2}} &&\\
B &= \dfrac{b^{2}\times b^{-9}}{b^{8}} &&\justif{$-3\times 3=-9$ et $4\times 2=8$}\\
B &= \dfrac{b^{-7}}{b^{8}} &&\justif{$2+(-9)=-7$}\\
B &= b^{-7-8} &&\\
B &= \res{b^{-15}} &&
\end{calculs}
\end{exo}

\partie{Notation scientifique}

\begin{exo}[Entiers et décimaux]
\explique{On place la virgule après le premier chiffre non nul, puis on compte de combien
de rangs elle s'est déplacée.}{elle recule : exposant positif ; elle avance : exposant
négatif}
\begin{calculs}
A &= 274~000 &&\justif{la virgule recule de $5$ rangs : $2{,}74$}\\
A &= \res{2{,}74\times 10^{5}} &&
\end{calculs}
\begin{calculs}
B &= 0{,}000~61 &&\justif{la virgule avance de $4$ rangs : $6{,}1$}\\
B &= \res{6{,}1\times 10^{-4}} &&
\end{calculs}
\begin{calculs}
C &= 38~200~000 &&\justif{la virgule recule de $7$ rangs}\\
C &= \res{3{,}82\times 10^{7}} &&
\end{calculs}
\begin{calculs}
D &= 0{,}000~000~29 &&\justif{la virgule avance de $7$ rangs}\\
D &= \res{2{,}9\times 10^{-7}} &&
\end{calculs}
\begin{calculs}
E &= 56~300 &&\justif{la virgule recule de $4$ rangs}\\
E &= \res{5{,}63\times 10^{4}} &&
\end{calculs}
\begin{calculs}
F &= 0{,}000~073~08 &&\justif{la virgule avance de $5$ rangs ; le $0$ entre $3$ et $8$ reste}\\
F &= \res{7{,}308\times 10^{-5}} &&
\end{calculs}
\rem{Contrôle rapide : un nombre plus grand que $10$ a un exposant positif, un nombre
compris entre $0$ et $1$ un exposant négatif.}
\end{exo}

\begin{exo}[Produits par une puissance de 10]
\explique{On écrit le premier facteur en notation scientifique, puis on regroupe les
puissances de $10$.}{$10^{n}\times 10^{p}=10^{n+p}$}
\begin{calculs}
A &= 532\times 10^{3} &&\justif{$532$ a trois chiffres avant la virgule : ce n'est pas fini}\\
A &= 5{,}32\times 10^{2}\times 10^{3} &&\justif{$532=5{,}32\times 10^{2}$}\\
A &= 5{,}32\times 10^{2+3} &&\\
A &= \res{5{,}32\times 10^{5}} &&
\end{calculs}
\begin{calculs}
B &= 64\times 10^{-5} &&\\
B &= 6{,}4\times 10^{1}\times 10^{-5} &&\justif{$64=6{,}4\times 10^{1}$}\\
B &= 6{,}4\times 10^{1+(-5)} &&\\
B &= \res{6{,}4\times 10^{-4}} &&
\end{calculs}
\begin{calculs}
C &= 481{,}7\times 10^{8} &&\\
C &= 4{,}817\times 10^{2}\times 10^{8} &&\justif{$481{,}7=4{,}817\times 10^{2}$}\\
C &= \res{4{,}817\times 10^{10}} &&\justif{$2+8=10$}
\end{calculs}
\begin{calculs}
D &= 0{,}036\times 10^{-4} &&\\
D &= 3{,}6\times 10^{-2}\times 10^{-4} &&\justif{$0{,}036=3{,}6\times 10^{-2}$ : la virgule avance de $2$ rangs}\\
D &= 3{,}6\times 10^{-2+(-4)} &&\\
D &= \res{3{,}6\times 10^{-6}} &&
\end{calculs}
\end{exo}

\begin{exo}[Un calcul jusqu'au bout]
\begin{calculs}
E &= \dfrac{35\times 10^{-2}\times 6\times(10^{3})^{3}}{14\times 10^{5}} &&\justif{on recopie}\\
E &= \dfrac{35\times 6}{14}\times\dfrac{10^{-2}\times 10^{9}}{10^{5}} &&\justif{les nombres d'un côté, les puissances de $10$ de l'autre}\\
E &= \dfrac{5\times\barre{7}\times\barre{2}\times 3}{\barre{2}\times\barre{7}}\times 10^{-2+9-5} &&\justif{$35=5\times 7$ ; $6=2\times 3$ ; $14=2\times 7$ ; $(10^{3})^{3}=10^{9}$}\\
E &= 15\times 10^{2} &&\justif{$5\times 3=15$ et $-2+9-5=2$}\\
E &= 1{,}5\times 10^{1}\times 10^{2} &&\justif{$15=1{,}5\times 10^{1}$}\\
E &= \res{1{,}5\times 10^{3}} &&
\end{calculs}
\rem{$15\times 10^{2}$ est juste mais n'est pas la notation scientifique : il faut un seul
chiffre non nul avant la virgule.}
\end{exo}

\begin{exo}[Deux expressions]
\cas{1)}
\begin{calculs}
A &= \dfrac{3\times 10^{1500}}{12\times 10^{1502}} &&\justif{on recopie}\\
A &= \dfrac{3}{12}\times\dfrac{10^{1500}}{10^{1502}} &&\justif{les nombres d'un côté, les puissances de $10$ de l'autre}\\
A &= \dfrac{1}{4}\times 10^{1500-1502} &&\justif{$\frac{3}{12}=\frac14$}\\
A &= 0{,}25\times 10^{-2} &&\justif{$\frac14=0{,}25$}\\
A &= 2{,}5\times 10^{-1}\times 10^{-2} &&\justif{$0{,}25=2{,}5\times 10^{-1}$}\\
A &= \res{2{,}5\times 10^{-3}} &&
\end{calculs}
\rem{Impossible d'écrire $10^{1500}$ en entier : ce sont les règles sur les puissances qui
font le calcul.}
\cas{2)}
\explique{Le numérateur de $B$ est une \emph{différence} : les règles sur les puissances
ne s'appliquent pas. On revient à l'écriture décimale.}{$10^{2}=100$ et
$10^{-1}=0{,}1$}
\begin{calculs}
B &= \dfrac{6{,}4\times 10^{2}-340\times 10^{-1}}{1{,}01\times 10^{2}} &&\\
B &= \dfrac{640-34}{101} &&\justif{$6{,}4\times 100=640$ ; $340\times 0{,}1=34$ ; $1{,}01\times 100=101$}\\
B &= \dfrac{606}{101} &&\\
B &= 6 &&\justif{$101\times 6=606$}
\end{calculs}
\explique{$B=6$ : \res{$B$ est un nombre entier}.}{}
\end{exo}

\begin{exo}[La lumière de la Lune]
\cas{1)}
\begin{calculs}
d &= 384~000\text{ km} &&\\
d &= 384~000~000\text{ m} &&\justif{$1$~km $=1~000$~m : trois zéros de plus}\\
d &= \res{3{,}84\times 10^{8}\text{ m}} &&\justif{la virgule recule de $8$ rangs}
\end{calculs}
\cas{2)}
\explique{En une seconde, la lumière parcourt $3\times 10^{8}$~m. Le temps, en secondes,
est le nombre de fois que cette longueur tient dans la distance Terre-Lune.}{temps $=$
distance $\div$ distance parcourue en une seconde}
\begin{calculs}
t &= \dfrac{3{,}84\times 10^{8}}{3\times 10^{8}} &&\justif{les deux distances en mètres}\\
t &= \dfrac{3{,}84}{3}\times\dfrac{10^{8}}{10^{8}} &&\justif{les nombres d'un côté, les puissances de $10$ de l'autre}\\
t &= 1{,}28\times 1 &&\justif{$\frac{10^{8}}{10^{8}}=10^{0}=1$}\\
t &= 1{,}28 &&
\end{calculs}
\explique{La lumière renvoyée par la Lune met environ \res{$1{,}28$~s} pour nous
parvenir.}{un peu plus d'une seconde}
\end{exo}

\end{document}
